% section: 等比数列の総和とその導出

  \section{等比数列の総和とその導出}
    発想の話をする.
    \[
      \sum_{k=1}^{n} a_1r^{n-1}=a+ar+ar^2+\cdots+ar^{n-1}
    \]
    これの求め方は
    \begin{align*}
      &S_{n}=a+ar+ar^2+\cdots+ar^{n-1} \tag*{\(\cdots\) (1)}\\
      &rS_n=ar+ar^2+ar^3+\cdots+ar^{n} \tag*{\(\cdots\) (2)}\\
    \end{align*}
    として,
    \[
    \begin{array}{rrrl}
        &S_{n}={}&a+&\hspace{-0.75em}ar+ar^2+\cdots+ar^{n-1}\\
      -\,) & rS_n={}&{}&\hspace{-0.75em}ar+ar^2+ar^3+\cdots+ar^{n}\\ \hline
        & (1-r)S_n={}&{}&\hspace{-0.75em}a-ar^n
    \end{array}
    \]
    \[
      \therefore \quad S_n=\frac{a(1-r^n)}{1-r}
    \]
    このように,数列においては総和を定数倍してから引くなどして,無限に存在しうる項を削除してしまう手法がよく取られる.
    覚えておこう.

